\documentclass[11pt]{article}
\usepackage[utf8]{inputenc}
\usepackage[T1]{fontenc}
\usepackage[margin=1in]{geometry}
\usepackage{amsmath,amssymb,amsthm}
\usepackage{booktabs}
\usepackage{array}
\usepackage[colorlinks=true,linkcolor=blue,citecolor=blue,urlcolor=blue]{hyperref}

\newtheorem{theorem}{Theorem}
\newtheorem{proposition}{Proposition}
\newtheorem{lemma}{Lemma}
\newtheorem{corollary}{Corollary}
\theoremstyle{definition}
\newtheorem{definition}{Definition}
\newtheorem{assumption}{Assumption}
\newtheorem{example}{Example}
\theoremstyle{remark}
\newtheorem*{remark}{Remark}

\newcommand{\Ev}{\mathcal{E}}
\newcommand{\K}{\mathcal{K}}
\newcommand{\Y}{\mathcal{Y}}
\newcommand{\Ex}{\mathbb{E}}
\newcommand{\R}{\mathbb{R}}
\newcommand{\tr}{\operatorname{tr}}
\newcommand{\rank}{\operatorname{rank}}
\newcommand{\dG}{d_{G}}
\newcommand{\PC}{P_{C}}

\title{\textbf{Geometric Evaluation Theory}\\[4pt]
\large Evaluator, Metric, and Budget Induce the Distinctions From Which Preference and Choice Follow}
\author{Andrew H. Bond\\
\small Department of Computer Engineering, San Jos\'e State University\\
\small \texttt{andrew.bond@sjsu.edu} ~\textbar~ ORCID: 0009-0003-2599-6158}
\date{September 2026. \emph{Draft 0.1.} Every statement is labeled proved, defined, or posited. Nothing labeled posited is a claim.}

\begin{document}
\maketitle

\begin{abstract}
We develop Geometric Evaluation Theory (GET), an evaluator-relative account in which preferences, values, and effective distinctions are induced by an evaluator's mapping into a consequence space, its metric of evaluation, and its finite resolution budget.
The primitive is not a utility ordering but an evaluation object, a six-tuple of states, actions, evaluator, metric, budget, and admissibility constraints.
From that object we derive the distinctions the evaluator can make, and from the distinctions we derive preference and choice.
At infinite budget the induced preference is a weak order with a utility, and the utility is the evaluator's distance to its ideal.
At finite resolution it is a semiorder with intransitive indifference, and the threshold is the budget rather than a free parameter.
Expected utility, mean--variance evaluation, additive multi-attribute value and its ideal-point ancestors in psychometrics and spatial voting, the geometric decision cost of our economics work, satisficing, and the decision geometry of Dawid and Lauritzen are recovered as special cases or degenerate projections, each with the conditions that make it one.
We then prove four things a scalar utility on the action set does not provide.
A change of resolution budget reverses preference between two fixed actions.
Two evaluators that agree on every pair of one menu can disagree on another, and they share a standard of correctness exactly when their induced distances are ordinally equivalent.
A representation shared by several evaluators at a finite budget is optimal for all of them exactly when their evaluation geometries share a leading eigenspace, and otherwise the least total regret is attained by the leading eigenspace of the summed geometries and is strictly positive.
Admissibility constraints that depend on what is available produce choice functions that violate the weak axiom of revealed preference, so an obligation of that kind is not a preference at any strength.
The theory claims a common mathematical grammar for economic, aesthetic, and moral evaluation and does not claim a common standard of correctness, and we say precisely where the difference between grammar and standard lies.
\end{abstract}

\section{Introduction}

Classical evaluation theory begins with an ordering.
An agent is assumed to rank alternatives, the ranking is assumed to satisfy axioms, and a utility function is constructed to represent it.
Everything that follows, expected utility, multi-attribute value, equilibrium, is a statement about that function.
The ordering itself is a primitive.

Geometric Evaluation Theory begins one step earlier.
An evaluator maps actions and states into a space of consequences.
It measures consequences with a metric of its own, from an ideal point of its own.
It does so at a finite resolution, so that some consequences it cannot tell apart.
The distinctions it can make are the first thing the theory derives, and preference and choice are derived from those distinctions.
The schematic is evaluator plus metric plus budget, then effective distinctions, then preference and choice.

This is a materially different foundation from the classical one, and it is different from the decision geometry of Dawid and Lauritzen~\cite{dawid2005,dawid2007}, which derives a geometry on the space of probability distributions from a loss function and its Bayes act.
That geometry appears here as the case of an evaluator that is the identity on distributions, an ideal that is the true distribution, and an infinite budget.
What the present theory adds is the evaluator map, which makes every geometry relative to who is evaluating, the budget, which makes every distinction relative to how finely they can look, and the admissibility constraints, which hold the part of evaluation that is not a trade-off.

The core of the induced preference, a weighted distance from an ideal point, is not new.
Coombs' unfolding models place alternatives and an ideal in a common space and rank by distance, and the spatial theory of voting of Davis, Hinich, and Ordeshook, developed by Enelow and Hinich, scores a policy by a weighted Euclidean distance from the voter's ideal, with a weight matrix that is exactly a constant metric of evaluation.
Those models are the unbudgeted, single-evaluator, identity-consumer case of the present object, and Section~\ref{sec:classical} says so.
What the present theory adds to them is the evaluator map, which places the ideal in a consequence space that the action space reaches only through a map the evaluator does not control, the budget, which turns the weak order of the ideal-point model into a semiorder and makes it depend on resolution, the admissibility set, and the multi-evaluator theorems.

The theory grew out of Observation Theory, where the same construction, a consumer map, an output metric, and a budget, was shown to induce the geometry that governs compression, retrieval, and estimation for a downstream reader.
An evaluator is an observer whose output is scored from an ideal point.
The theorems of that program transfer, and we cite them rather than reprove them.

The essential discipline is to claim a common mathematical grammar and not a common standard of correctness.
An economic preference, an aesthetic judgment, and a moral obligation may share the same geometric form without being reducible to the same scalar utility.
We make that discipline precise in two places.
Two evaluators share a standard exactly when their induced distances are ordinally equivalent, which the grammar does not impose.
And admissibility constraints that depend on the menu produce choice that no preference ordering rationalizes, so obligations of that kind are outside preference altogether rather than preferences with a large weight.

The paper does five things.
Section~\ref{sec:object} defines the evaluation object and separates what the world supplies from what the evaluator owns.
Section~\ref{sec:derive} derives distinctions, preference, and choice from it and states the representation problem that the separation creates.
Section~\ref{sec:classical} recovers the classical models with their conditions.
Section~\ref{sec:new} proves the four results that scalar evaluation does not provide.
Section~\ref{sec:predict} states what the theory predicts, in a form a measurement can fail.
Section~\ref{sec:branches} says how the descriptive, normative, and multi-agent branches of the program sit under the theory, with what is measured and what is posited in each.

\section{The Evaluation Object}
\label{sec:object}

\begin{definition}[Evaluation object]
\label{def:object}
An evaluation object is a six-tuple
\[
\Ev=(X,A,C,G,B,\K)
\]
in which
\begin{enumerate}
\item[(i)] $X$ is a set of states carrying a probability law $\mu$, the evaluator's belief, which may be a point mass when there is no uncertainty,
\item[(ii)] $A$ is a set of actions or alternatives,
\item[(iii)] $C:A\times X\to\Y$ is the evaluator, a map into a consequence space $\Y$ that is a $C^{1}$ manifold, with $C(a,\cdot)$ measurable for each $a$,
\item[(iv)] $G$ is the metric of evaluation, a pointed geometry on $\Y$ consisting of a continuous Riemannian metric, also written $G$, and an ideal point $t\in\Y$,
\item[(v)] $B$ is the resolution budget, one of a rank cap $k$, a length $\varepsilon\ge0$, or a rate $R$ in bits, with $B=\infty$ for the unbudgeted case,
\item[(vi)] $\K\subseteq A$ is the set of admissible actions, which may depend on the state and on the menu presented.
\end{enumerate}
\end{definition}

Three remarks on the definition.
The ideal point belongs to the metric because a metric alone orders nothing.
A distance becomes an evaluation when it is measured from somewhere, and where an evaluator measures from is as much a part of its standard as how it measures.
The evaluator is a map, not a function of the action alone, because consequences depend on the state and the evaluator's belief about the state is part of the object.
The admissible set is a set and not a cost, because the theory needs a place for evaluation that is not a trade-off, and Section~\ref{sec:new} shows that place cannot be a term in a cost.

\begin{definition}[World and evaluator]
\label{def:split}
The world supplies $(X,A,C)$, the states, the actions, and the map from actions and states to consequences.
The evaluator supplies $(G,B,\K)$, its metric with its ideal, its budget, and its admissible set.
Two evaluators of the same world differ only in the second triple.
\end{definition}

The split is the foundational commitment of the theory and the source of its content.
If the evaluator could also choose $C$, every weak order on a finite $A$ would be an evaluation object, since a one-dimensional consequence $C(a)=t-d(a)$ represents any real function $d$, and the theory would say nothing.
With $C$ fixed by the world, the question of which preferences an evaluator can hold becomes a representation problem, stated at the end of Section~\ref{sec:derive}.

\begin{definition}[Evaluation cost and distance]
\label{def:cost}
The evaluation distance of an action is the expected geodesic distance of its consequence from the ideal,
\[
d(a)=\Bigl(\Ex_{\mu}\bigl[\dG\bigl(C(a,X),t\bigr)^{2}\bigr]\Bigr)^{1/2},
\]
and its evaluation cost is $c(a)=d(a)^{2}$.
When $X$ is a point, $d(a)=\dG(C(a),t)$.
\end{definition}

\begin{definition}[Budgeted distance]
\label{def:budget}
Let $\Y=\R^{m}$ with $G$ constant, so that $\dG(y,t)=\|G^{1/2}(y-t)\|$, and let $\omega$ be the evaluator's workload, a probability law on $A\times X$ that describes what it habitually evaluates.
Let $M_{\omega}=\Ex_{\omega}\bigl[G^{1/2}(C-t)(C-t)^{\top}G^{1/2}\bigr]$ be the second moment of evaluated consequences.
\begin{enumerate}
\item[(i)] Under a rank cap $k$, let $\Pi_{k}$ be the orthogonal projector onto a top-$k$ eigenspace of $M_{\omega}$. The budgeted distance is $d_{k}(a)=\bigl(\Ex_{\mu}\|\Pi_{k}G^{1/2}(C(a,X)-t)\|^{2}\bigr)^{1/2}$.
\item[(ii)] Under a length $\varepsilon$, the budgeted distance is $d$ itself and the budget acts on comparison, in Definition~\ref{def:pref} below.
\item[(iii)] Under a rate $R$, the evaluator holds a description of the consequence at $R$ bits and evaluates the reconstruction. The consumer-relative rate--distortion function of the observation-theory paper gives the least distortion $D(R)$ attainable, and we treat a rate budget as the length budget $\varepsilon(R)=D(R)^{1/2}$.
\end{enumerate}
On a general manifold the same definitions hold locally through the pullback expansion of the observation-theory representation theorem.
\end{definition}

\begin{remark}
The workload $\omega$ is exogenous to any particular menu.
An evaluator that instead selects its retained directions from the menu in front of it is the budgeted menu-relative observer of the observation-theory paper, and Section~\ref{sec:new} recalls what that observer does to regularity.
\end{remark}

\section{What the Object Induces}
\label{sec:derive}

\subsection{Effective distinctions}

\begin{theorem}[Induced distinctions]
\label{thm:dist}
Let $\Ev$ be an evaluation object.
\begin{enumerate}
\item[(a)] \emph{Unbudgeted.} The relation $a\sim b$ if and only if $C(a,x)=C(b,x)$ for $\mu$-almost every $x$ is an equivalence relation on $A$. When $A$ is a manifold and $C$ is $C^{1}$ in $a$, the pullback $\PC(a)=J_{C}(a)^{\top}G(C(a))J_{C}(a)$ is positive semidefinite with kernel $\ker J_{C}(a)$, and at regular points the classes are locally submanifolds tangent to that kernel.
\item[(b)] \emph{Rank cap.} The relation $a\sim_{k}b$ if and only if $\Pi_{k}G^{1/2}C(a,x)=\Pi_{k}G^{1/2}C(b,x)$ for $\mu$-almost every $x$ is an equivalence relation, each of whose classes is a union of classes of $\sim$, and $\sim_{k'}$ refines $\sim_{k}$ when $k'>k$ and the eigenspaces are nested.
\item[(c)] \emph{Length.} The relation $a\approx_{\varepsilon}b$ if and only if $|d(a)-d(b)|\le\varepsilon$ is reflexive and symmetric. It is transitive on $A$ if and only if $\varepsilon=0$ or the set $\{d(a):a\in A\}$ contains no three points $r_{1}<r_{2}<r_{3}$ with $r_{2}-r_{1}\le\varepsilon$, $r_{3}-r_{2}\le\varepsilon$, and $r_{3}-r_{1}>\varepsilon$.
\end{enumerate}
\end{theorem}

\begin{proof}
(a) is Theorem 1(a) and 1(b) of the observation-theory representation theorem~\cite{bond2026or} applied to the consumer $C(\cdot,x)$ at each state and intersected over $\mu$-almost every state, since equality almost everywhere is an equivalence and the kernel of an integral of positive semidefinite forms is the almost-everywhere common kernel.
(b) The relation is equality of a function of $a$, hence an equivalence, and $C(a,x)=C(b,x)$ implies $\Pi_{k}G^{1/2}C(a,x)=\Pi_{k}G^{1/2}C(b,x)$, so each $\sim_{k}$ class contains the $\sim$ classes of its members. If the top-$k$ eigenspace is contained in the top-$k'$ eigenspace then equality after $\Pi_{k'}$ implies equality after $\Pi_{k}$.
(c) Reflexivity and symmetry are immediate. If three such points exist then $r_{1}\approx_{\varepsilon}r_{2}$ and $r_{2}\approx_{\varepsilon}r_{3}$ but not $r_{1}\approx_{\varepsilon}r_{3}$. If none exist, take $a\approx_{\varepsilon}b\approx_{\varepsilon}c$ with values $r_{a},r_{b},r_{c}$. If $r_{a}\le r_{b}\le r_{c}$ or the reverse, the absence of such a triple gives $|r_{a}-r_{c}|\le\varepsilon$, and in every other ordering $|r_{a}-r_{c}|\le\max(|r_{a}-r_{b}|,|r_{b}-r_{c}|)\le\varepsilon$.
\end{proof}

Part (c) is the whole difference between a resolution budget and a rank budget.
A rank budget coarsens the evaluator's world into classes.
A length budget does not produce classes at all.
It produces a tolerance, and the failure of transitivity is not a defect of the evaluator but the signature of finite resolution.

\subsection{Preference}

\begin{definition}[Induced preference]
\label{def:pref}
Fix $\Ev$ and write $d_{B}$ for the budgeted distance, with $d_{\infty}=d$ and $d_{\varepsilon}=d$.
Strict preference and indifference are
\[
a\succ_{B}b \iff d_{B}(a)+\varepsilon_{B}<d_{B}(b),
\qquad
a\sim_{B}b \iff |d_{B}(a)-d_{B}(b)|\le\varepsilon_{B},
\]
where $\varepsilon_{B}=\varepsilon$ for a length budget, $\varepsilon_{B}=\varepsilon(R)$ for a rate budget, and $\varepsilon_{B}=0$ otherwise.
\end{definition}

\begin{theorem}[Preference from geometry]
\label{thm:pref}
Let $A$ be finite.
\begin{enumerate}
\item[(a)] For $\varepsilon_{B}=0$ the relation $\succsim_{B}$, defined as $\succ_{B}$ or $\sim_{B}$, is a weak order represented by the utility $u_{B}(a)=-d_{B}(a)$, and $u_{B}$ is unique up to strictly increasing transformation among representations of $\succsim_{B}$.
\item[(b)] For $\varepsilon_{B}>0$ the pair $(\succ_{B},\sim_{B})$ is a semiorder in the sense of Luce, and $u_{B}=-d_{B}$ with threshold $\varepsilon_{B}$ is a threshold representation of it, $a\succ_{B}b$ if and only if $u_{B}(a)>u_{B}(b)+\varepsilon_{B}$.
\item[(c)] The semiorder of (b) is a weak order if and only if the condition of Theorem~\ref{thm:dist}(c) holds, and in particular as $\varepsilon_{B}\to0$ it converges to the weak order of (a).
\end{enumerate}
\end{theorem}

\begin{proof}
(a) A relation defined by comparing a real function is complete and transitive, and the function represents it. Uniqueness up to monotone transformation is the standard ordinal uniqueness.
(b) A relation defined by $u(a)>u(b)+\varepsilon$ with fixed $\varepsilon>0$ is irreflexive and satisfies the two semiorder axioms. If $a\succ b$ and $c\succ e$, then $u(a)>u(b)+\varepsilon$ and $u(c)>u(e)+\varepsilon$, so either $u(a)>u(e)+\varepsilon$ or $u(c)>u(b)+\varepsilon$, since if both failed we would have $u(a)\le u(e)+\varepsilon<u(c)$ and $u(c)\le u(b)+\varepsilon<u(a)$, a contradiction. If $a\succ b\succ c$ and $e$ is any action, then $u(a)>u(c)+2\varepsilon$, so either $u(a)>u(e)+\varepsilon$ or $u(e)\ge u(a)-\varepsilon>u(c)+\varepsilon$. This is the easy direction of the Scott and Suppes representation of semiorders, and the representation is by construction.
(c) The semiorder is a weak order exactly when indifference is transitive, which is Theorem~\ref{thm:dist}(c), and for finite $A$ the condition holds for all sufficiently small $\varepsilon$.
\end{proof}

\begin{remark}
Luce~\cite{luce1956} derived semiorders from just-noticeable differences in utility with the threshold as a primitive, Scott and Suppes~\cite{scott1958} gave the representation, and Fishburn~\cite{fishburn1970} and Tversky~\cite{tversky1969} the extensions and the evidence.
Here the threshold is the budget and the utility is the pullback distance to the ideal, so both are fixed by the evaluation object.
A semiorder with intransitive indifference is therefore what a finite-resolution evaluator must exhibit, and its threshold is a measurable property of the evaluator rather than a fitted constant.
\end{remark}

\subsection{Choice}

\begin{definition}[Choice]
\label{def:choice}
For a menu $M\subseteq A$, the choice set is the set of admissible actions that no admissible action in the menu strictly beats,
\[
\mathrm{ch}_{\Ev}(M)=\bigl\{a\in M\cap\K(M):\ \text{there is no } b\in M\cap\K(M)\text{ with } b\succ_{B}a\bigr\}.
\]
A population choice frequency, where one is wanted, is $p(a\mid M)\propto\exp(-c_{B}(a)/T)$ over $M\cap\K(M)$ with a temperature $T$, and it is a modeling choice outside the derivation.
\end{definition}

\begin{proposition}[Choice is nonempty and satisficing is a budget]
\label{prop:choice}
For finite $M$ with $M\cap\K(M)$ nonempty, $\mathrm{ch}_{\Ev}(M)$ is nonempty.
Under a length budget, every admissible action within $\varepsilon$ of the least admissible distance is chosen, and if the ideal is attainable to within $\varepsilon$ then every admissible action within $\varepsilon$ of the ideal is chosen.
\end{proposition}

\begin{proof}
A semiorder on a finite set is acyclic, so it has a maximal element. The second statement is Definition~\ref{def:pref} read at the minimum.
\end{proof}

The second statement is Simon's satisficing~\cite{simon1955}, with the aspiration level as the ideal point and the tolerance as the budget.
An evaluator that stops at the first good-enough action is not departing from the geometry.
It is the geometry at finite resolution.

\subsection{The representation problem}

\begin{definition}[Representable preference]
Fix a world $(X,A,C)$ with $\Y=\R^{m}$.
A weak order $\succsim$ on $A$ is representable if there is a constant positive semidefinite $G$ and an ideal $t$ with $a\succsim b$ if and only if $d(a)\le d(b)$ for the distance of Definition~\ref{def:cost}.
A semiorder is representable at budget $\varepsilon$ if the same holds with Definition~\ref{def:pref}.
\end{definition}

\begin{proposition}[Two facts about representability]
\label{prop:repr}
\begin{enumerate}
\item[(a)] If $X$ is a point and the consequences $C(A)$ are affinely independent in $\R^{m}$, then every weak order on $A$ with at most $m+1$ elements is representable, and the representation is not unique.
\item[(b)] If $X$ is a point and $C(A)$ lies on a line in $\R^{m}$, then a weak order on $A$ is representable if and only if it is single-peaked along that line, with the peak at the point of the line nearest the ideal.
\end{enumerate}
\end{proposition}

\begin{proof}
(a) Choose $G=I$. The map $t\mapsto(\|y_{a}-t\|^{2})_{a\in A}$ is, up to the common term $\|t\|^{2}$, affine in $t$ with coefficient vectors $-2y_{a}$ and constants $\|y_{a}\|^{2}$. Affine independence of at most $m+1$ points makes the differences of these affine functions linearly independent, so any prescribed ordering of the values, including ties, is attained by some $t$, and an open set of $t$ attains any strict ordering.
(b) Along a line with coordinate $s$, the distance to any ideal is $\sqrt{(s-s_{0})^{2}g+h}$ for the projection $s_{0}$ of the ideal onto the line and constants $g>0$, $h\ge0$, which is decreasing then increasing in $s$. Conversely every single-peaked order is attained by placing the ideal's projection at the peak with a suitable $g$.
\end{proof}

Part (b) is the spatial-voting single-peakedness condition, and it shows the content of the split.
When the world confines consequences to a line, the evaluator cannot hold a preference with two peaks, whatever its metric and ideal.
The general characterization of representable preferences for a fixed world of dimension $m$ and more than $m+1$ actions is the foundational open problem of the theory, and we do not claim it.

\section{Classical Models as Special Cases}
\label{sec:classical}

Each row of Table~\ref{tab:classical} names an evaluation object, and each is proved in the proposition that follows.

\begin{table}[h]
\centering
\small
\begin{tabular}{p{0.24\textwidth} p{0.50\textwidth} p{0.20\textwidth}}
\toprule
Model & Evaluation object & What is degenerate \\
\midrule
Expected utility & $\Y=\R$, $C(a)=\Ex_{\mu}v(o(a,X))$ for a von Neumann--Morgenstern index $v$, $G=1$, $t\ge\sup C$, $B=\infty$, $\K=A$ & one dimension, no budget \\
Mean--variance & $\Y=\R$, $C(a,x)=v(o(a,x))$, $G=1$, $t\ge\sup v$, $B=\infty$ & uncertainty inside the metric \\
Additive multi-attribute value & $\Y=\R^{m}$, $G=\operatorname{diag}(w)$, $t$ the ideal profile, $B=\infty$ & second-order form near $t$ \\
Geometric decision cost & $\Y=\R^{9}$, $C$ the scenario encoding, $G=\Sigma^{-1}$, $t$ the task reference, $B=$ rank $3$ as sparsity of $G$, $\K$ the menu grid & the descriptive branch \\
Satisficing & any $\Y$, $t$ the aspiration, $B=\varepsilon>0$ & the length budget alone \\
Dawid--Lauritzen decision geometry & $\Y$ the probability simplex, $C=\mathrm{id}$, $G=\nabla^{2}(-H)$ for the generalized entropy $H$ of a proper scoring rule, $t$ the true distribution, $B=\infty$ & identity evaluator, no budget \\
Nash equilibrium & $N$ objects with $C_{i}$ affine in own action, $G_{i}=H_{i}$, $t_{i}$ the own-action maximizer & the multi-agent branch \\
\bottomrule
\end{tabular}
\caption{Classical models as evaluation objects.}
\label{tab:classical}
\end{table}

\begin{proposition}[Recoveries]
\label{prop:classical}
\begin{enumerate}
\item[(a)] With the expected-utility object, $a\succsim b$ if and only if $\Ex v(o(a,X))\ge\Ex v(o(b,X))$.
\item[(b)] With the mean--variance object, $c(a)=(t-\Ex v)^{2}+\operatorname{Var}v$, so preference is decreasing in the mean and in the variance of $v(o(a,X))$, and the object is the expected-utility object with the expectation taken after rather than before the metric.
\item[(c)] With the additive multi-attribute object, $c(a)=\sum_{i}w_{i}(y_{i}(a)-t_{i})^{2}$, whose gradient at a profile $y$ is $2w_{i}(y_{i}-t_{i})$, so the additive value function with weights $2w_{i}(t_{i}-y_{i})$ is its first-order form at $y$.
\item[(d)] The geometric decision cost object reproduces Corollary 2 of the observer-representation paper~\cite{bond2026or}, and the sixteen reported predictions of the economics paper~\cite{bond2026tcss}, by construction.
\item[(e)] The satisficing object yields Proposition~\ref{prop:choice}.
\item[(f)] For a proper scoring rule with concave generalized entropy $H$ that is $C^{2}$ on the interior of the simplex, the associated divergence $D(P,Q)$ satisfies $D(P,P+\delta)=\tfrac12\delta^{\top}\nabla^{2}(-H)(P)\,\delta+o(\|\delta\|^{2})$, so the Dawid--Lauritzen decision geometry is the evaluation cost of the stated object to second order.
\item[(g)] The multi-agent object is the behavioral game of the observer-representation paper, whose equilibria coincide with Nash equilibria for utilities of the form $K_{i}(s_{-i})-c_{i}(a_{i};s_{-i})$ and for concave quadratic games.
\end{enumerate}
\end{proposition}

\begin{proof}
(a) $c(a)=(t-\Ex v)^{2}$ is decreasing in $\Ex v$ on $\Ex v\le t$.
(b) Expand $\Ex[(t-v)^{2}]$ about the mean.
(c) Differentiate.
(d) The object is the one written in that corollary.
(e) Immediate.
(f) The divergence of a proper scoring rule is the Bregman divergence of the convex function $-H$, and a Bregman divergence agrees with the quadratic form of the Hessian to second order.
(g) This is Theorem 4 of the observer-representation paper.
\end{proof}

\begin{remark}
The ideal-point models of Coombs~\cite{coombs1964} and of spatial voting theory~\cite{davis1970,enelow1984} are the additive multi-attribute row with $X$ a point and $C$ the identity on a policy space, so they are in the table under that row.
Two classical objects are absent from the table because the theory as written does not contain them.
Loss aversion is an asymmetry between directions of the same displacement, and a Riemannian metric is symmetric, so reference-dependent evaluation with loss aversion needs the metric of evaluation to be a Finsler metric, which we have not developed.
A lexicographic priority between attributes is the limit of $G=\operatorname{diag}(1,\delta)$ as $\delta\to0$, but at $\delta=0$ it is not a metric, and by Debreu's theorem~\cite{debreu1954} it is not a utility either.
Section~\ref{sec:new} places priorities of that kind in $\K$.
\end{remark}

\section{What Scalar Evaluation Does Not Provide}
\label{sec:new}

Each result below is a theorem about evaluation objects and a statement about what a utility function on $A$ cannot do.

\subsection{Resolution reverses preference}

\begin{theorem}[Resolution reversal]
\label{thm:reversal}
There is an evaluation object and two actions $a,b$ such that $b\succ_{1}a$ under rank cap $1$ and $a\succ_{2}b$ under rank cap $2$.
Consequently no single function $u:A\to\R$ represents the evaluator's preference at both budgets, while the family $u_{k}=-d_{k}$ does.
\end{theorem}

\begin{proof}
Take $\Y=\R^{2}$, $G=I$, $t=0$, $X$ a point, and a workload whose consequence second moment has top eigenvector $e_{1}$, for example $M_{\omega}=\operatorname{diag}(4,1)$.
Let $C(a)=(1,0)$ and $C(b)=(0.9,2)$.
Then $d_{1}(a)=1$, $d_{1}(b)=0.9$, $d_{2}(a)=1$, and $d_{2}(b)=\sqrt{0.81+4}=2.19$.
A function representing both would satisfy $u(b)>u(a)$ and $u(a)>u(b)$.
\end{proof}

\begin{remark}
The reversal is not a change of taste.
The evaluator, its metric, and its ideal are the same at both budgets.
What changed is how many directions of consequence it can afford to resolve, and an action that is good in the direction it resolves first and bad in the direction it resolves second is preferred at the coarse budget and rejected at the fine one.
A scalar utility has no parameter in which to record that.
\end{remark}

\subsection{Evaluator-dependent distinctions and the standard of correctness}

\begin{theorem}[Grammar without standard]
\label{thm:standard}
Let $\Ev_{1}$ and $\Ev_{2}$ be unbudgeted evaluation objects on the same actions $A$ with induced distances $d_{1},d_{2}$.
\begin{enumerate}
\item[(a)] Their induced equivalence relations coincide if and only if $C_{1}(a,\cdot)=C_{1}(b,\cdot)$ almost everywhere exactly when $C_{2}(a,\cdot)=C_{2}(b,\cdot)$ almost everywhere. For $C^{1}$ evaluators on a manifold $A$ this holds locally at regular points if and only if $\ker J_{C_{1}}=\ker J_{C_{2}}$.
\item[(b)] They induce the same preference on every menu if and only if there is a strictly increasing $\varphi$ with $d_{2}=\varphi\circ d_{1}$ on $A$.
\item[(c)] Agreement on every pair of one menu $M$ does not imply agreement on another. There are objects that agree on a menu of any finite size and disagree on a menu that adds one action.
\end{enumerate}
\end{theorem}

\begin{proof}
(a) is the definition of the relation, and the local statement is Theorem~\ref{thm:dist}(a).
(b) Preference is the ordering of $d_{i}$ by Theorem~\ref{thm:pref}(a), and two real functions on a set induce the same weak order if and only if one is a strictly increasing transform of the other on that set.
(c) Take $\Y=\R^{2}$, $G_{1}=\operatorname{diag}(1,\delta)$, $G_{2}=\operatorname{diag}(\delta,1)$, $t=0$, and a menu whose actions all have consequences on the diagonal $y_{1}=y_{2}$, on which $d_{1}$ and $d_{2}$ coincide up to the common factor $\sqrt{1+\delta}$. Add an action with consequence $(1,0)$. Under $G_{1}$ it has distance $1$, under $G_{2}$ it has distance $\sqrt{\delta}$, and for small $\delta$ it is ranked last by the first evaluator and first by the second among actions on the diagonal at distance between $\sqrt{\delta}$ and $1$.
\end{proof}

Part (b) is where the difference between grammar and standard lies.
Two evaluation objects share the grammar by construction.
They share a standard of correctness exactly when their induced distances are ordinally equivalent, and nothing in the six-tuple imposes that.
An economic evaluator, an aesthetic evaluator, and a moral evaluator can each be written as an object, and writing them so asserts nothing about whether they agree.

\subsection{Coupled evaluators and the impossibility of a neutral representation}

\begin{theorem}[Incompatibility of coupled geometries]
\label{thm:incompat}
Let $N$ evaluators share a representation space $\R^{d}$ in which consequences are evaluated, with constant pullbacks $P_{1},\dots,P_{N}$ and a common isotropic source of second moment $\sigma^{2}I$ on that space.
Let a shared representation at budget $k$ be an orthogonal projector $\Pi$ of rank $k$, reconstructing $\Pi x$.
Evaluator $i$'s regret from $\Pi$ is
\[
r_{i}(\Pi)=\sigma^{2}\Bigl(\sum_{j\le k}\lambda_{j}(P_{i})-\tr(\Pi P_{i})\Bigr)\ge0 .
\]
\begin{enumerate}
\item[(a)] A shared representation with $r_{i}(\Pi)=0$ for every $i$ exists if and only if there is a $k$-dimensional subspace that is a top-$k$ eigenspace of every $P_{i}$.
\item[(b)] For weights $w_{i}>0$, the least weighted total regret $\min_{\Pi}\sum_{i}w_{i}r_{i}(\Pi)$ is attained by a top-$k$ eigenspace of $\bar{P}=\sum_{i}w_{i}P_{i}$ and equals $\sigma^{2}\bigl(\sum_{i}w_{i}\sum_{j\le k}\lambda_{j}(P_{i})-\sum_{j\le k}\lambda_{j}(\bar{P})\bigr)$, which is strictly positive exactly when (a) fails.
\end{enumerate}
\end{theorem}

\begin{proof}
The regret formula is Proposition 2 of the observer-representation paper~\cite{bond2026or} applied to each evaluator, and both steps below are Ky Fan's maximum principle~\cite{fan1949}.
(a) $r_{i}(\Pi)=0$ if and only if $\Pi$ projects onto a top-$k$ eigenspace of $P_{i}$, by the equality case of Ky Fan's maximum principle, and a common $\Pi$ is a common such subspace.
(b) $\sum_{i}w_{i}\tr(\Pi P_{i})=\tr(\Pi\bar{P})$ is maximized over rank-$k$ projectors on a top-$k$ eigenspace of $\bar{P}$ with value $\sum_{j\le k}\lambda_{j}(\bar{P})$, again by Ky Fan, which gives the minimum and its value. The value is zero if and only if every term is zero, since each $r_{i}\ge0$, which is (a).
\end{proof}

\begin{remark}
The theorem is about institutions, shared codes, shared ontologies, and shared training objectives, anything several evaluators must read through the same finite representation.
It says that a representation neutral among evaluators exists only when their evaluation geometries share a leading subspace, and that otherwise the best shared representation is the leading subspace of the influence-weighted sum of their geometries, with a positive regret that the theorem computes.
No scalar utility contains the statement, because the statement is about subspaces and not about levels.
\end{remark}

\subsection{Admissibility is not a preference}

\begin{theorem}[Menu-dependent admissibility violates revealed preference]
\label{thm:admiss}
Let $\K(M)$ depend on the menu, so that an action admissible in one menu can be inadmissible in a larger one, and let choice be Definition~\ref{def:choice}.
Then there are evaluation objects and menus $M\subset M'$ with $a,b\in M$ such that $a$ is chosen from $M$, $b$ is not chosen from $M$, and $b$ is chosen from $M'$.
The choice function violates the weak axiom of revealed preference, and no preference ordering on $A$, and hence no utility, rationalizes it.
\end{theorem}

\begin{proof}
Let $A=\{a,b,p\}$, $d(a)<d(b)<d(p)$, $\K(\{a,b\})=\{a,b\}$, and $\K(\{a,b,p\})=\{b,p\}$, as when the availability of $p$ places $a$ under an obligation that removes it.
Then $\mathrm{ch}(\{a,b\})=\{a\}$ and $\mathrm{ch}(\{a,b,p\})=\{b\}$.
The weak axiom requires that if $b$ is chosen from a menu containing $a$ then $a$ is not chosen from any menu containing $b$ from which $b$ is not chosen, and it fails.
A choice function rationalized by a preference ordering satisfies the axiom.
The menus are recorded as presented, with $a$ available in both, since the admissibility is the evaluator's and not the world's.
\end{proof}

\begin{remark}
A fixed admissible set is representable by a utility with a penalty, so the theorem is not about constraints as such.
It is about constraints that depend on what else is available, which is the form obligations take when promises, consent, rights created by others' actions, or roles are in play.
That menu-dependent norms defeat revealed preference is Sen's observation~\cite{sen1971,sen1993}, and the present theorem is his argument written for evaluation objects.
What the theory adds is a place for such constraints in the object beside the evaluator, the metric, and the budget, so that a moral obligation and an economic preference can be written in one grammar while the obligation stays outside the cost.
This is the formal content of the claim that the theory does not reduce morality to preference maximization.
\end{remark}

\subsection{The menu-relative budgeted evaluator}

\begin{proposition}[Regularity]
\label{prop:reg}
An evaluator whose rank budget is applied to the second moment of the menu in front of it rather than to an exogenous workload has a choice function that violates regularity and therefore lies outside every random-utility model.
\end{proposition}

\begin{proof}
This is the non-reducibility theorem of the observer-representation paper, with the explicit menu given there.
\end{proof}

\section{What the Theory Predicts}
\label{sec:predict}

A foundation that only recovers known models is a notation.
The following are the theory's predictions, each stated so that a measurement can fail it, and none is measured yet.
\begin{enumerate}
\item[(i)] \emph{Threshold tracks budget.} For a fixed world and evaluator, the indifference threshold of the induced semiorder is the resolution budget. A manipulation that reduces resolution, time pressure, load, or a coarser display of consequences, must enlarge the just-noticeable difference of Theorem~\ref{thm:pref}(b) without changing the ordering of well-separated pairs. A change in the ordering of well-separated pairs under such a manipulation fails the prediction.
\item[(ii)] \emph{Rank reversal.} A manipulation that changes which consequence directions the evaluator resolves, without changing its metric or ideal, must produce the reversals of Theorem~\ref{thm:reversal} between actions that trade the resolved direction against an unresolved one, and must not produce reversals between actions that differ only in the resolved directions.
\item[(iii)] \emph{Shared-representation regret.} When several evaluators read one representation of a fixed budget, the loss each suffers relative to its own best representation is the regret of Theorem~\ref{thm:incompat}, and the representation that minimizes the influence-weighted total is the leading eigenspace of the weighted sum of their geometries. A shared representation that does better than that bound for every evaluator fails the prediction.
\item[(iv)] \emph{Menu-dependent admissibility.} Where an obligation is created by the availability of an action, choice must violate the weak axiom in the pattern of Theorem~\ref{thm:admiss}, and the violation must disappear when the obligation-creating action is removed from every menu.
\end{enumerate}
The first two require an evaluator whose budget can be manipulated and whose consequence map is known, which is a designed experiment.
The third requires several evaluators whose geometries can be recovered, which is the readscope instrument applied to each.
The fourth requires menus with and without an obligation-creating action, which is a choice experiment.

\section{The Three Branches}
\label{sec:branches}

The theory sits above three branches of the program, and above Observation Theory, which is the theory of an evaluator whose ideal is the true value of what it reads.
For each branch we say what the evaluation object is, what has been measured, and what is posited.

\paragraph{Descriptive choice.}
The economics paper is the object of Table~\ref{tab:classical}, row four.
Measured are the sixteen aggregate targets, the rejection of stake invariance, and the nonmonotone temperature that carries the ordering of two of the lottery targets.
Posited are the nine coordinates and every constant of the encoding.
Nothing in that paper tests the budget as a property of the reader, and the learned-geometry protocol of the observer-representation paper is the test.

\paragraph{Normative judgment.}
The ethics stack of the program~\cite{bond2026encyclopedia} writes a moral evaluation as a signed valence on nine frozen coordinates, the moral dimensions, with a deontic gate that admits or refuses an action before any valence is weighed.
In the present vocabulary the coordinates are $\Y$, the valence and its reliability weights are $G$, the gate is $\K$, and the Bond Invariance Principle, that a judgment is valid only if invariant under bond-preserving transformations, is the requirement that the evaluator factor through the quotient by the group of such transformations, which is the statement that the group's orbits lie in the fibers of Theorem~\ref{thm:dist}(a).
Measured are the encoders' held-out discriminations on eight of nine coordinates, with one reported failed, and the shakedown-level consumer-weighted invariance results.
Posited are the coordinates themselves and the correspondence just stated, which the program's registry lists as not yet canonized.
Aesthetic judgment is an evaluation object with no action set, an evaluator on works, a metric, and an ideal, and the program's measurement there is the compressibility result of the aesthetics volume, which we do not restate.

\paragraph{Multi-agent interaction.}
The behavioral game of the observer-representation paper is $N$ evaluation objects with an exogenous target for each, and Theorem~\ref{thm:incompat} is its institutional face.
Proved are existence, the Nash correspondence on the stated class, and hardness.
Posited is every application to a real institution.

\section{What Is Proved and What Is Not}

\begin{center}
\small
\begin{tabular}{p{0.52\textwidth} p{0.12\textwidth} p{0.30\textwidth}}
\toprule
Statement & Status & Where \\
\midrule
The evaluation object and its budgets & defined & Definitions~\ref{def:object} to \ref{def:budget} \\
Unbudgeted distinctions are an equivalence; rank budgets coarsen it; length budgets give a tolerance that is transitive only under a stated condition & proved & Theorem~\ref{thm:dist} \\
Preference is a weak order with utility $-d$ at zero threshold and a semiorder with threshold equal to the budget otherwise & proved & Theorem~\ref{thm:pref} \\
Choice is nonempty; satisficing is the length budget & proved & Proposition~\ref{prop:choice} \\
Expected utility, mean--variance, additive value, decision cost, satisficing, Dawid--Lauritzen geometry, and Nash are special cases & proved with conditions & Proposition~\ref{prop:classical} \\
Loss aversion and lexicographic priority are within the theory & not claimed & remark after Proposition~\ref{prop:classical} \\
A rank budget change reverses preference between fixed actions & proved & Theorem~\ref{thm:reversal} \\
Shared grammar, shared standard iff ordinal equivalence, agreement on a menu does not transfer & proved & Theorem~\ref{thm:standard} \\
Neutral shared representation iff common leading eigenspace; least regret at the leading eigenspace of the weighted sum & proved & Theorem~\ref{thm:incompat} \\
Menu-dependent admissibility violates the weak axiom & proved & Theorem~\ref{thm:admiss} \\
Representability for affinely independent consequences and single-peakedness on a line & proved & Proposition~\ref{prop:repr} \\
The general representation theorem for a fixed world & open, not claimed & Section~\ref{sec:derive} \\
The four predictions & stated, unmeasured & Section~\ref{sec:predict} \\
The ethics-stack correspondence, gate as $\K$ and invariance as factoring through the quotient & posited & Section~\ref{sec:branches} \\
Any claim that a human evaluator has a measured budget & not claimed & Section~\ref{sec:branches} \\
\bottomrule
\end{tabular}
\end{center}

\begin{thebibliography}{99}

\bibitem{dawid2005}
A.~P.~Dawid and S.~L.~Lauritzen, ``The geometry of decision theory,'' in \emph{Proc. 2nd Int. Symp. Information Geometry and Its Applications}, Univ. of Tokyo, 2005, pp.~22--28.

\bibitem{dawid2007}
A.~P.~Dawid, ``The geometry of proper scoring rules,'' \emph{Ann. Inst. Statist. Math.}, vol.~59, no.~1, pp.~77--93, 2007, doi: 10.1007/s10463-006-0099-8.

\bibitem{luce1956}
R.~D.~Luce, ``Semiorders and a theory of utility discrimination,'' \emph{Econometrica}, vol.~24, no.~2, pp.~178--191, 1956, doi: 10.2307/1905751.

\bibitem{scott1958}
D.~Scott and P.~Suppes, ``Foundational aspects of theories of measurement,'' \emph{J. Symbolic Logic}, vol.~23, no.~2, pp.~113--128, 1958, doi: 10.2307/2964389.

\bibitem{fishburn1970}
P.~C.~Fishburn, ``Intransitive indifference with unequal indifference intervals,'' \emph{J. Math. Psychol.}, vol.~7, no.~1, pp.~144--149, 1970, doi: 10.1016/0022-2496(70)90062-3.

\bibitem{tversky1969}
A.~Tversky, ``Intransitivity of preferences,'' \emph{Psychol. Rev.}, vol.~76, no.~1, pp.~31--48, 1969, doi: 10.1037/h0026750.

\bibitem{simon1955}
H.~A.~Simon, ``A behavioral model of rational choice,'' \emph{Quart. J. Econ.}, vol.~69, no.~1, pp.~99--118, 1955, doi: 10.2307/1884852.

\bibitem{debreu1954}
G.~Debreu, ``Representation of a preference ordering by a numerical function,'' in \emph{Decision Processes}, R.~M.~Thrall, C.~H.~Coombs, and R.~L.~Davis, Eds. New York, NY, USA: Wiley, 1954, pp.~159--165.

\bibitem{sen1971}
A.~K.~Sen, ``Choice functions and revealed preference,'' \emph{Rev. Econ. Stud.}, vol.~38, no.~3, pp.~307--317, 1971, doi: 10.2307/2296384.

\bibitem{sen1993}
A.~Sen, ``Internal consistency of choice,'' \emph{Econometrica}, vol.~61, no.~3, pp.~495--521, 1993, doi: 10.2307/2951715.

\bibitem{coombs1964}
C.~H.~Coombs, \emph{A Theory of Data}. New York, NY, USA: Wiley, 1964.

\bibitem{davis1970}
O.~A.~Davis, M.~J.~Hinich, and P.~C.~Ordeshook, ``An expository development of a mathematical model of the electoral process,'' \emph{Amer. Polit. Sci. Rev.}, vol.~64, no.~2, pp.~426--448, 1970, doi: 10.2307/1953842.

\bibitem{enelow1984}
J.~M.~Enelow and M.~J.~Hinich, \emph{The Spatial Theory of Voting: An Introduction}. Cambridge, U.K.: Cambridge Univ. Press, 1984.

\bibitem{keeney1993}
R.~L.~Keeney and H.~Raiffa, \emph{Decisions with Multiple Objectives: Preferences and Value Trade-Offs}. Cambridge, U.K.: Cambridge Univ. Press, 1993.

\bibitem{kahneman1979}
D.~Kahneman and A.~Tversky, ``Prospect theory: An analysis of decision under risk,'' \emph{Econometrica}, vol.~47, no.~2, pp.~263--291, 1979, doi: 10.2307/1914185.

\bibitem{fan1949}
K.~Fan, ``On a theorem of Weyl concerning eigenvalues of linear transformations I,'' \emph{Proc. Nat. Acad. Sci. USA}, vol.~35, no.~11, pp.~652--655, 1949, doi: 10.1073/pnas.35.11.652.

\bibitem{bond2026or}
A.~H.~Bond, ``Observer representation: One pullback, two pushforwards,'' draft 0.1, 2026, in the geometric-observation repository, doi: 10.5281/zenodo.21776291.

\bibitem{bond2026tcss}
A.~H.~Bond, ``Geometric prediction of economic behavior: Cross-domain validation across game theory and prospect theory,'' \emph{IEEE Trans. Comput. Social Syst.}, 2026, accepted for publication.

\bibitem{bond2026encyclopedia}
A.~H.~Bond, \emph{The Observation Theory Encyclopedia}, 2026. [Online]. Available: \url{https://erisml.org/encyclopedia}

\end{thebibliography}

\end{document}
